\par\bigskip
\pdfbookmark[1]{Bemerkungen}{purae-bemerkungen}
\begin{center}\textsc{Bemerkungen.}\end{center}
\begingroup\small
Von der ursprünglichen Fortsetzung dieser Abhandlung von art. 19 an, welche der Behandlung specieller Fälle gewidmet war, sind nur noch einige Artikel vorhanden, die sich mit der quadratischen Gleichung beschäftigen, deren Wurzeln die beiden $\frac{n-1}2$-gliedrigen Perioden sind; das Manuscript bricht im Anfang der Untersuchung ab, durch welche das Vorzeichen der bei der Auflösung derselben auftretenden Quadratwurzel bestimmt werden sollte; aus der Uebereinstimmung dieses noch vorhandenen Anfangs mit der Abhandlung Summatio quarumdam serierum singularium geht hervor, dass der Verfasser seinen Plan änderte, um die eben erwähnte Bestimmung des Vorzeichens zum Gegenstande einer besondern Abhandlung zu machen. Vergleicht man hiermit das Citat im art. 8 (wo im Manuscript statt der zweiten Ausgabe des Werkes von \textsc{Lagrange} durch ein Versehen die dritte angegeben war), so ergibt sich, dass diese Handschrift aus dem Jahre 1808 stammt. Dass aber die Publication des Vorhergehenden nicht aufgegeben war, lehrt ein bei art. 19 offenbar in späterer Zeit eingeschobenes Blatt, auf welchem eine andere Fortsetzung beginnt und bezüglich der Bestimmung des Vorzeichens schon auf die Abhandlung Summatio etc. verwiesen wird. Diese zweite Fortsetzung, welche aber auch bald abbricht, ist hier mitgetheilt. Der Text des durchaus druckfertigen Manuscriptes ist bei der Herausgabe treu beibehalten, nur in art. 16 mussten die Formeln für den zweiten Fall hinzugefügt werden.
\begin{flushright}R. \textsc{Dedekind}.\end{flushright}
\endgroup
\BookRule{.34\textwidth}
